\section{Кольцо}
\label{sec:Ring}
\tikzsetfigurename{LinearAlgebra_Ring_}

\begin{definition}[Кольцо]
    Непустое множество $R$ с двумя бинарными операциями $\opAdd: R \times R \to R$ (сложение) и $\opMult: R \times R \to R$ (умножение) называется \emph{кольцом}, если выполнены следующие условия.
    \begin{enumerate}
        \item $(R, \opAdd)$~--- это абелева группа, нейтральный элемент которой~--- $\Bbbzero$.
              Для любых $a, b, c \in R$:
              \begin{itemize}
                  \item $(a \opAdd b) \opAdd c = a \opAdd (b \opAdd c)$
                  \item $\Bbbzero \opAdd a = a \opAdd \Bbbzero = a$
                  \item $a \opAdd b = b \opAdd a$
                  \item для любого $a \in R$ существует $-a \in  R$, такой что $a + (-a) = \Bbbzero$.
              \end{itemize}
              В последнем пункте кроется отличие от полукольца.
              \item $(R, \opMult)$~--- это моноид, нейтральный элемент которого~--- $\Bbbone$.
              Для любых $a, b, c \in R$:
              \begin{itemize}
                  \item $(a \opMult b) \opMult c = a \opMult (b \opMult c)$
                  \item $\Bbbone \opMult a = a \opMult \Bbbone = a$
              \end{itemize}
        \item $\opMult$ дистрибутивно слева и справа относительно $\opAdd$:
              \begin{itemize}
                  \item $a \opMult (b \opAdd c) = (a \opMult b) \opAdd (a \opMult c)$
                  \item $(a \opAdd b) \opMult c = (a \opMult c) \opAdd (b \opMult c)$
              \end{itemize}
    \end{enumerate}
\end{definition}

Заметим, что мультипликативное свойство $\Bbbzero$ (быть аннигилятором по умножению) не указывается явно, так как может быть выведено из остальных утверждений.
Действительно,
\begin{enumerate}
    \item $a \opMult \Bbbzero = a \opMult (\Bbbzero \opAdd \Bbbzero)$, так как $\Bbbzero$~--- нейтральный по сложению, то $\Bbbzero \opAdd \Bbbzero = \Bbbzero$
    \item Воспользуемся дистрибутивностью: $a \opMult (\Bbbzero \opAdd \Bbbzero) = a \opMult \Bbbzero \opAdd a \opMult \Bbbzero$.
          В итоге: $a \opMult \Bbbzero = a \opMult \Bbbzero \opAdd a \opMult \Bbbzero$
    \item Так как у нас есть группа по сложению, то для любого $a$ существует обратный элемент $a^{-1}$, $a \opAdd a^{-1} = \Bbbzero$.
          Прибавим $a^{-1} \opMult \Bbbzero$ к левой и правой части равенства%
          \sidenote{Обычно данное действие воспринимается как очевидное, но, строго говоря, оно требует аккуратного введения структур с равенством и соответствующих аксиом.}%
          , полученного на предыдущем шаге:
          \[a \opMult \Bbbzero \opAdd a^{-1} \opMult \Bbbzero = a \opMult \Bbbzero \opAdd a \opMult \Bbbzero \opAdd a^{-1} \opMult \Bbbzero.\]
    \item Воспользуемся дистрибутивностью и ассоциативностью.
          \begin{align*}
              (a \opAdd a^{-1}) \opMult \Bbbzero & = a \opMult \Bbbzero \opAdd (a  \opAdd a^{-1}) \opMult \Bbbzero \\
              \Bbbzero \opMult \Bbbzero          & = a \opMult \Bbbzero \opAdd \Bbbzero \opMult \Bbbzero           \\
              \Bbbzero                           & = a \opMult \Bbbzero
          \end{align*}
    \item Аналогично можно доказать, что $\Bbbzero = \Bbbzero \opMult a$.
\end{enumerate}

%\section{Поле}
